\subsection{Fermionic hypercharge transport on prepared source addresses}
\label{sec:fermion-source-current}

Fix the same 64 prepared golden addresses and 144 oriented nearest-neighbour
links as the scalar execution. Give each address two Spin components of
each of the five registered matter multiplets. Their internal multiplicities
and integer charges are
\[
 (d_s)=(6,3,3,2,1),\qquad (q_s)=6(Y_s)=(1,-4,2,-3,6).
\]
This finite model supplies its spatial hopping action, canonical
anticommutation relations, state preparation and model step. The source
addresses do not select those inputs. The colour and weak indices are
retained as multiplicities of an abelian transport law; the model does not
execute the complete nonabelian or Yukawa dynamics.

Let $c_{s,a,v,\alpha}$ be the annihilation operators on the finite fermion
Fock space, where $a=1,\ldots,d_s$ and $\alpha=1,2$. They satisfy
$\{c_i,c_j^\dagger\}=\delta_{ij}$ and $\{c_i,c_j\}=0$.
For an oriented link $e=(u,v)$ in direction $j$, a unit complex link $U_e$
and a supplied real hopping weight $\kappa_e$, put
\begin{align}
 K_{s,e}&=i\sigma_j U_e^{q_s},\\
 \widehat H_e&=\kappa_e\sum_{s,a}
 \bigl(c_{s,a,u}^\dagger K_{s,e}c_{s,a,v}
       +c_{s,a,v}^\dagger K_{s,e}^\dagger c_{s,a,u}\bigr).
 \label{eq:fermion-source-edge-action}
\end{align}
The Pauli matrices act on Spin, not on the internal multiplicity index.
The sign and link orientation in this definition fix the current convention.
If $U_e\mapsto g_uU_e\overline g_v$ and
$c_{s,a,w}\mapsto g_w^{q_s}c_{s,a,w}$, the quadratic expression is
unchanged. The link is held fixed during this matter update.

Occupy one normalized spatial--Spin orbital $\psi_s$ for every internal
index $a$. Distinct internal indices and multiplets are orthogonal, so their
exterior product is a normalized fifteen-fermion Slater state. Its one-body
density is the orthogonal projector
\[
 P=\bigoplus_s I_{d_s}\otimes |\psi_s\rangle\langle\psi_s|.
\]
The complex orbital coefficients describe this fermionic state. They are
not commuting replacements for the field operators. In particular
$0\leq P\leq1$ and $\operatorname{Tr}P=15$.

For one occupied orbital write $z=(\psi_{s,u},\psi_{s,v})$ and
\[
 h_{s,e}=\kappa_e
 \begin{pmatrix}0&K_{s,e}\\K_{s,e}^\dagger&0\end{pmatrix}.
\]
For step $\delta>0$, take the implicit midpoint equation
\begin{equation}
 i\frac{z'-z}{\delta}=h_{s,e}\bar z,
 \qquad \bar z=\frac{z'+z}{2}.
 \label{eq:fermion-source-midpoint}
\end{equation}
Since $K_{s,e}$ is unitary, $h_{s,e}^2=\kappa_e^2 I$.
The update is the unitary Cayley transform
$(I+i\delta h_{s,e}/2)^{-1}(I-i\delta h_{s,e}/2)$.
It preserves the projector property and particle number, with
$P\mapsto VPV^\dagger$, and is exact over Gaussian rationals
when the link, step, weight and preparation lie in that field.

\begin{proposition}[Variational current and mean Gauss preservation]
\label{prop:fermion-source-current}
Define the mean site charge and midpoint current by
\begin{align}
 \rho_w&=\sum_s d_s q_s\|\psi_{s,w}\|^2,\\
 J_e&=-2\kappa_e\sum_s d_sq_s\,
       \operatorname{Im}(\bar\psi_{s,u}^\dagger
                         K_{s,e}\bar\psi_{s,v}).
 \label{eq:fermion-source-current}
\end{align}
Then $J_e$ is the derivative of the quadratic midpoint functional of
\eqref{eq:fermion-source-edge-action} under $U_e\mapsto e^{i\theta}U_e$,
with the midpoint orbitals fixed. These averages need not be normalized
occupied orbitals; only the step endpoints represent the normalized
fifteen-fermion Slater states. The charge changes satisfy
\[
 \rho'_u-\rho_u=-\delta J_e,\qquad
 \rho'_v-\rho_v=\delta J_e.
\]
Let $E_e$ be a classical link momentum and update
$E'_e=E_e-\delta J_e$. With $\operatorname{div}$ positive at the initial
endpoint of each oriented link, every site residual
$G_w=(\operatorname{div}E)_w-\rho_w$ is unchanged. This remains true for
any finite composition of these edge factors.
\end{proposition}
\begin{proof}
For the finite canonical anticommutation relations,
$[c_i^\dagger c_j,c_k^\dagger c_l]
=\delta_{jk}c_i^\dagger c_l-\delta_{li}c_k^\dagger c_j$;
the degree-four terms cancel with their fermionic signs. Thus the quadratic
Hamiltonian acts through its one-body matrix. Differentiating its link
phase at fixed state gives the displayed $J_e$.
The identity
\[
 \|x'\|^2-\|x\|^2
 =2\operatorname{Re}\left(\frac{x'+x}{2}\right)^\dagger(x'-x)
\]
and \eqref{eq:fermion-source-midpoint} give opposite endpoint transfers.
The electric update cancels those changes in $G_u,G_v$; all other site
charges and incident momenta are unchanged. Induction proves the finite
composition statement. Gauge covariance follows by conjugating the
one-body matrices and their Cayley transforms.
\end{proof}

Write $U_e=e^{i\theta_e}$. The same first-order action is
\[
 S_e=\int\left(i\sum_s d_s\psi_s^\dagger\dot\psi_s
       +E_e\dot\theta_e-\langle\widehat H_e\rangle\right)dt.
\]
Its variations give the orbital equation, $\dot E_e=-\partial_{\theta_e}
\langle\widehat H_e\rangle$ and $\dot\theta_e=0$, since this factor has no
$E_e$ dependence in its Hamiltonian. Thus frozen links and electric feedback
belong to one declared action factor. The midpoint scheme discretizes
these equations. Exact discrete conservation does not identify its step
with the continuous Hamiltonian exponential or prove a trajectory-error
estimate. The electric momentum is driven by a fermionic expectation:
this is a semiclassical update, not a quantized gauge field coupled to
fermion Fock space. Electric drift, plaquette and Yukawa factors are absent.

\paragraph{A sharp quantum-constraint boundary.}
At every address choose $\psi_s=(3/40,1/10)$ before the local phase
intervention. Each orbital has squared site norm $1/64$ and unit global
norm. Since $\sum_s d_sq_s=0$, the initial mean charge vanishes at every
site, and $E=0$ solves the mean Gauss constraint. Nevertheless the local
charge operator has variance
\begin{equation}
 \operatorname{Var}(\widehat Q_w)
 =\sum_s d_sq_s^2\frac1{64}\left(1-\frac1{64}\right)
 =\frac{945}{512}>0.
 \label{eq:fermion-source-gauss-variance}
\end{equation}
Each occupied internal orbital contributes an independent Bernoulli site
occupation; this proves the formula. A site phase preserves it. Consequently
the product preparation with sharp classical electric data is not
annihilated by the quantum Gauss operator. In ordinary hypercharge units
the variance is divided by $36$. Vanishing mean residuals cannot certify
a physical gauge-constrained quantum state.

\paragraph{Execution and interpretation.}
The executable construction in \path{code/sm_fermion_current/} retains
the local orbital state, incident link ports, current readback and electric
feedback as versioned software records. Its independent verifier rebuilds
the charge dictionary, source addresses, Cayley updates, variational
currents and consumed writer/value versions. Baseline, local phase
intervention and gauge-related execution are separate controls. The
generation's linear, cubic and mixed anomaly sums are representation
identities; a finite canonical-anticommutation implementation does not
construct a chiral continuum fermion measure.

This packet has one declared generation and supplied quantum preparation.
It proves neither native source selection, all-matter Standard Model
evolution, a controlled spatial or temporal continuum limit, laboratory
current, a physical clock, nor compatibility with a sourced curved family.
The preceding scalar reduction and its continuous-readout enclosure retain
their own action and preparation; their error bounds do not transfer to
this fermionic transport.

\subsection{A flux-dressed fermion state satisfying abelian quantum Gauss}
\label{sec:fermion-quantum-link}

The mean-field obstruction has a precise remedy within a separately
supplied quantum link model. It changes both the gauge-field Hilbert space
and the preparation. Choose one oriented link $(u,v)$ of the same golden
carrier. Its electric space is $\ell^2(\mathbb Z)$, with
\[
 \widehat E|n\rangle=n|n\rangle,\qquad
 \widehat U|n\rangle=|n+1\rangle,\qquad
 [\widehat E,\widehat U]=\widehat U.
\]
The ambient flux space is infinite. A finite reachable support is computed
from the preparation; no cyclic wrap or truncated shift replaces this
unitary operator.

Index the fifteen internal channels by $a$, retaining their charges $q_a$.
In each channel occupy the normalized left Spin orbital
$\chi=(3/5,4/5)$ and leave its orthogonal orbital empty. The active right
orbital is $-i\sigma_j\chi$, so the Pauli link maps it back to $\chi$.
The two active modes have annihilators $l_a,r_a$. Fix the canonical
fermion order $(l_1,r_1,l_2,r_2,\ldots)$ and set
\begin{equation}
 \widehat h_a=\kappa_a
 \left(l_a^\dagger\widehat U^{q_a}r_a
       +r_a^\dagger\widehat U^{-q_a}l_a\right).
 \label{eq:quantum-link-fermion-action}
\end{equation}
Each channel retains exactly one fermion. Distinct channel Hamiltonians
commute: their even fermion bilinears have disjoint mode support and the
flux shifts commute. Initially all fifteen fermions are at $u$, the right
modes are empty, and the link is in $|0\rangle$. Since
$\sum_a q_a=0$, this preparation is locally neutral.

\begin{proposition}[Exact finite-support abelian quantum Gauss state]
\label{prop:fermion-quantum-link-gauss}
Let
$\widehat Q_u=\sum_aq_a l_a^\dagger l_a$ and
$\widehat Q_v=\sum_aq_a r_a^\dagger r_a$.
Every $\widehat h_a$ commutes with
\[
 \widehat G_u=\widehat E-\widehat Q_u,\qquad
 \widehat G_v=-\widehat E-\widehat Q_v.
\]
Its Cayley transform preserves their joint kernel. Starting from the
neutral preparation, the product of these transforms is the normalized
state
\begin{equation}
 |\Psi\rangle=
 \sum_{x\in\{0,1\}^{15}}
 \left[\prod_a A_a^{1-x_a}(-iB_a)^{x_a}\right]
 |x\rangle_{\!F}\otimes
 \left|-\sum_aq_ax_a\right\rangle_{\!E},
 \label{eq:fermion-quantum-link-state}
\end{equation}
where $x_a=1$ means that channel $a$ occupies its right mode, and
\[
 \eta_a=\delta\kappa_a/2,\qquad
 A_a=\frac{1-\eta_a^2}{1+\eta_a^2},\qquad
 B_a=\frac{2\eta_a}{1+\eta_a^2}.
\]
Every branch, and hence the complete state, satisfies both quantum Gauss
constraints. All other vertices retain zero charge and incident flux.
\end{proposition}
\begin{proof}
In $l_a^\dagger\widehat U^{q_a}r_a$, both left charge and electric
flux increase by $q_a$, while right charge decreases by $q_a$. The
adjoint reverses all three changes. This proves the two commutators on
the invariant finite-support core of the electric operator. The
Hamiltonians are bounded; their Cayley transforms preserve the joint Gauss
sectors and the electric domain. On the one-fermion channel sector,
$\widehat h_a^2=\kappa_a^2I$, which gives the displayed two coefficients.
The interleaved fermion convention fixes the even-hop signs. Channel
commutation gives the product expansion. Its orthogonal Fock branches
have squared norms summing to
$\prod_a(A_a^2+B_a^2)=1$.
On branch $x$, $Q_v=\sum_aq_ax_a$, $Q_u=-Q_v$ and $E=-Q_v$.
Both Gauss residuals therefore vanish individually; no cancellation of
opposite residuals in an expectation is used.
\end{proof}

The span of these $2^{15}$ dressed branches is invariant under every
$\widehat h_a$: each admissible hop changes the occupation and flux by the
same signed charge. Thus the finite electric support is exact for the
coupled evolution of this sector, although the standalone shift
$\widehat U$ does not preserve it. Channel commutation also gives the
continuous solution of the supplied Hamiltonian $\sum_a\widehat h_a$:
replace $A_a,B_a$ in \eqref{eq:fermion-quantum-link-state} by
$\cos(\kappa_at),\sin(\kappa_at)$. This is an exact finite-sector quantum
history, not a time-step approximation.

For $\delta=1/2$, choose $\kappa_a=1$ except on the $e^c$ channel,
where $\kappa_a=2$. These species weights are declared; they do not
establish a universal physical propagation speed. The coefficients are
$(A,B)=(15/17,8/17)$ on the first fourteen channels and $(3/5,4/5)$ on
the last. All $2^{15}$ Fock branches occur, with electric support
$\{-18,\ldots,18\}$. The link-phase current operator is
\[
 \widehat J_a=iq_a\kappa_a
 \left(l_a^\dagger\widehat U^{q_a}r_a
       -r_a^\dagger\widehat U^{-q_a}l_a\right),
 \qquad \dot{\widehat E}=-\sum_a\widehat J_a.
\]
Exact evaluation in \eqref{eq:fermion-quantum-link-state} gives
\[
 \langle\widehat E\rangle=-\frac{18144}{7225},\qquad
 \sum_a\langle\widehat J_a\rangle=\frac{47232}{7225}.
\]
The sum of the factor-midpoint currents is instead $36288/7225$;
multiplied by $-\delta$, it equals the electric change. A final-state
current and an integrated discrete transfer are different readouts.

For the continuous solution with the same species weights,
\[
 \langle\widehat E\rangle_t=6(\sin^2t-\sin^2 2t),\qquad
 \langle\widehat J\rangle_t=-6\sin 2t+12\sin 4t.
\]
The electric derivative equals minus the current. At the supplied model
time $t=\pi/6$ they give $-3$ and $3\sqrt3$, respectively, with electric
variance $45/2$. The flux fluctuations are correlated with charge, so quantum Gauss
holds throughout this history. The time parameter, action and preparation
are inputs; this solution calibrates no physical clock.

The executable producer and independent verifier in
\path{code/sm_fermion_current/quantum_link.py} and
\path{code/sm_fermion_current/verify_quantum_link.py} bind the channel and
Spin maps, CAR signs, complete branch population, flux distribution and
current matrix elements. This is a supplied finite quantum state and
action factor on one prepared link. It is not the uniform semiclassical
Slater preparation, a measured quantum history, full gauge-field
evolution, nonabelian Gauss constraints or Yukawa evolution, or a chiral continuum
regulator.
Native state and operation selection, source-causal attachment and
laboratory realization remain additional requirements.
